\section{Introducción}

La agricultura representa uno de los pilares económicos de Colombia, contribuyendo significativamente al Producto Interno Bruto (PIB) y al empleo rural. Sin embargo, en regiones como el departamento del Huila, caracterizado por su producción de café, cacao, frutas y hortalizas, los productores enfrentan barreras estructurales que limitan su competitividad y rentabilidad. Una de las problemáticas más acuciantes es la intermediación excesiva en la cadena de comercialización, donde múltiples actores (comisionistas, mayoristas y distribuidores) capturan una porción desproporcionada del valor agregado, dejando a los agricultores con márgenes reducidos y vulnerables a fluctuaciones de precios \cite{fao2020intermediation}.

El proyecto "Portal Agro Comercial del Huila" emerge como una iniciativa innovadora para mitigar estos desafíos. Este portal digital busca integrar tecnologías de la información y comunicación (TIC) para crear un marketplace virtual que conecte directamente a productores locales con compradores nacionales e internacionales, facilitando transacciones transparentes, trazabilidad de productos y acceso a información de mercado en tiempo real. Al reducir la dependencia de intermediarios tradicionales, el portal no solo optimiza la comercialización sino que también promueve prácticas sostenibles, fortalece la resiliencia de las cadenas de suministro y contribuye al desarrollo económico regional.

Este artículo académico proporciona un análisis profundo y estructurado del marco teórico que sustenta la implementación del portal, basado en 19 artículos clave identificados por su relevancia en la adopción tecnológica en el agro colombiano y latinoamericano. Cada artículo se analiza en detalle, incluyendo explicaciones profundas, aplicaciones prácticas al proyecto, potenciales mejoras e innovaciones implementables, y sugerencias de visualizaciones gráficas. Además, se incorpora una revisión exhaustiva de literatura adicional reciente (2020-2025), obtenida de fuentes como Google Scholar y bases académicas especializadas, para ampliar y actualizar el análisis.

La metodología adoptada combina un enfoque cualitativo de revisión sistemática de literatura con un análisis comparativo entre autores y perspectivas. Se identifican patrones comunes en la adopción tecnológica, barreras específicas del contexto colombiano, y oportunidades emergentes derivadas de tendencias globales como la Agricultura 4.0 y el comercio electrónico rural.

El artículo se organiza de la siguiente manera: la Sección 2 presenta el marco teórico, con subsecciones dedicadas a cada uno de los 19 artículos base. La Sección 3 detalla los hallazgos de la literatura adicional. La Sección 4 propone aplicaciones específicas al proyecto Portal Agro Comercial del Huila. La Sección 5 describe un modelo conceptual del portal. La Sección 6 sugiere gráficas y visualizaciones para ilustrar los conceptos clave. Finalmente, la Sección 7 ofrece conclusiones orientadas a la reducción de intermediación en el Huila, destacando implicaciones para políticas públicas, oportunidades de inversión y recomendaciones para escalabilidad.

Este trabajo contribuye al campo de la agricultura digital en contextos en desarrollo, proporcionando una base teórica sólida para iniciativas similares en Colombia y Latinoamérica. Al integrar evidencia empírica con propuestas prácticas, se busca no solo validar la viabilidad del portal sino también inspirar su adopción como modelo de transformación digital en el sector agropecuario.